\documentclass[12pt]{article}
\usepackage[utf8]{inputenc}
\usepackage[T1]{fontenc}
\usepackage{hyperref}
\title{Solano Language Resources for K-12 Curriculum Development}
\author{Compiled by Claude Sonnet 5.5 (claude-sonnet-5-5) for Andr\'as Kornai}
\date{}
\begin{document}
\maketitle

\section{Introduction}
Solano (ISO 639-3 xso; Glottocode sanf1266) is an extinct, unclassified language of northeastern Mexico, and possibly southern Texas. It is known only from a 21-word vocabulary list at the end of a 1703--1708 baptism book from the San Francisco Solano Mission, which served several peoples. It has been grouped with the proposed Coahuiltecan languages, a grouping most linguists now reject. Wikipedia dates its extinction to the 18th century.

\section{Primary Educational Resources}
None exist. No primers, textbooks or teaching materials were found, and with a 21-word record none could be built from the language itself.

\section{Linguistic and Reference Materials}
\begin{itemize}
\item Wikipedia, ``Solano language''. \url{https://en.wikipedia.org/wiki/Solano_language}
\item Wikipedia, ``Solano people''. \url{https://en.wikipedia.org/wiki/Solano_people}
\item Wikipedia, ``Coahuiltecan languages'', for the proposed grouping. \url{https://en.wikipedia.org/wiki/Coahuiltecan_languages}
\item UBC Wiki, ``Linguistic material from the tribes of Southern Texas and Northeastern Mexico''. I did not check what it contains for Solano. \url{https://wiki.ubc.ca/Documentation:RelLex/Linguistic_material_from_the_tribes_of_Southern_Texas_and_Northeastern_Mexico}
\item INAH journal article that appeared in the search results; I did not verify its relevance to Solano. \url{https://revistas.inah.gob.mx/index.php/noroestedemexico/article/download/21844/23230/45540}
\end{itemize}

\section{Religious and Community Resources}
No community organisations or religious materials in Solano were found. The only source is the mission baptism book mentioned above.

\section{Digital and Audio Resources}
No audio, video or software was found.

\section{Curriculum Development Context}
A language curriculum in Solano is not feasible on the evidence found. Any heritage programme would have to be about the history of the Solano people and the San Francisco Solano Mission, and the regional Indigenous context, rather than the language. The 21-word list could be mentioned as a historical document. Consult descendant communities in Coahuila and Texas, if they exist, before any such use.
\end{document}